% ============================================================
%  CHƯƠNG 7: KẾT LUẬN VÀ HƯỚNG PHÁT TRIỂN
%  Người viết: [Họ tên thành viên phụ trách]
% ============================================================

\section{Kết luận}
\label{sec:conclusion}

% -----------------------------------------------------------
\subsection{Tóm tắt kết quả đạt được}
% -----------------------------------------------------------
% Gợi ý nội dung:
%   - Tóm lại những gì đã trình bày: lý thuyết GMM, suy dẫn EM, thực nghiệm
%   - Kết quả nổi bật từ phần thực nghiệm
%   - Trả lời câu hỏi: GMM-EM có ưu điểm gì so với k-means?

% TODO: Viết nội dung ở đây


% -----------------------------------------------------------
\subsection{Hạn chế}
% -----------------------------------------------------------
% Gợi ý nội dung:
%   - Hội tụ cục bộ (local optimum), phụ thuộc khởi tạo
%   - Cần xác định trước K
%   - Chi phí tính toán khi D và K lớn
%   - Vấn đề singularity của ma trận hiệp phương sai

% TODO: Viết nội dung ở đây


% -----------------------------------------------------------
\subsection{Hướng phát triển}
% -----------------------------------------------------------
% Gợi ý nội dung:
%   - Variational EM và Bayesian GMM (tự động xác định K)
%   - Dirichlet Process Mixture Model (vô hạn thành phần)
%   - Cải tiến tốc độ: mini-batch EM, online EM
%   - Ứng dụng mở rộng: Hidden Markov Model, VAE

% TODO: Viết nội dung ở đây
